\chapter{Conclusion}

In this project, we have priced American put options on both the Black-Scholes and Heston models with various techniques, building benchmarks in the process to evaluate the validity of our results. The project mainly focuses on the use of variational methods to solve the variational inequality associated with the American option pricing problem, with analysis on the numerical results for their strengths and shortcomings. This project builds upon pre-existing techniques such as the tree and the PINN method, and adapts them to suit the American option pricing problem, while also providing a comprehensive mathematical justification for these numerical methods. We go beyond the existing literature by applying the variational methods to the problem of finding the free boundary where we develop a novel soft classification scheme, as well as constructing a solid foundation in our benchmarking process as a sanity check for our results.

We can see that the variational methods are very versatile: they are simply a matter of changing the underlying PDE and modifying the existing boundary conditions to apply the same method to a different model. Compare this with the tree method we used for the benchmark: the addition of a stochastic process for the variance in the Heston model requires a disproportionate amount of extra effort and care to construct a recombining model. The variational methods are also more efficient than the tree method: once trained, the neural network can be evaluated at any point in the domain very efficiently as it is only a forward pass, while the tree method requires the construction of the entire tree.

The versatility of the variational methods is underpinned by the exploration of the behaviour of the free boundary of the put on two assets after changing the payoffs. Even though we do not have a valid benchmark (a tree method would be difficult to construct, and a closed-form solution is of course not available), we can see that the free boundary behaves as expected. Granted, our model sacrifices some accuracy as we scale up the complexity and the number of dimensions, but the results are still reasonable, also taking into consideration hardware constraints.

A drawback of the variational methods is that they are not as accurate as the tree method, since they rely on approximating the solution with only the governing PDE and boundary conditions to guide the training. We can observe systematic bias in the results in different parts of the domain space, leading to less desirable results in the free boundary approximation with the Gaussian assumption, where the unbiasedness assumption is violated. Additionally, the variational methods are not as interpretable as the tree method, which has a clear structure in its discretisation of the underlying stochastic process and its pricing architecture is more transparent with its backwards induction. As a result, the variational methods are much more difficult to debug and can actually be relatively unstable to train, as tuning various hyperparameters and changing the activation function can lead to very different results. This brings into question the feasibility of this method in the practical world, where the lack of interpretability and stability can be a problem for risk management, as we do not, of course, have any benchmarks to evaluate the results against.

For future work, it would be interesting to explore yet more complex models: we have shown that the variational method is versatile enough to be applied to any model with any payoff, so we could introduce either a different volatility process or add a stochastic interest rate. Additionally, the work on the multi-dimensional pricing problem could be further refined, where we take advantage of a GPU to train a larger and more capable neural network across a longer training time, which would hopefully lead to better results. As lattice methods suffer from the curse of dimensionality, we could consider using the Least-Squares Monte Carlo method as a benchmark, which is more efficient in higher dimensions. Finally, the work on the free boundary approximation could be further refined, where we work on better soft classification models and more complex Monte Carlo schemes to retrieve the continuation value, which would hopefully allow us to observe some interesting shapes in complex payoffs in higher dimensions.
